% !PRINTED 120
\weberpage{120}
%<<CONT

\[
 2y=u+v+w
\]
und erhalten, wenn wir zur Abkürzung
\begin{equation}
 s=u^2+v^2+w^2,\qquad
 t=v^2w^2+w^2u^2+u^2v^2
\tag{2}
\end{equation}
setzen,
\[
 4y^2=s+2(vw+wu+uv),
\]
\[
 16y^4=s^2+4s(vw+wu+uv)+4t+8uvw(u+v+w).
\]
Wenn man dies in (1) einsetzt, so findet sich
\[
 s^2+4t+4as+16c+8(uvw+b)(u+v+w)
  +4(s+2a)(vw+wu+uv)=0.
\]
Diese Gleichung wird aber durch die Annahme befriedigt
\[
 s+2a=0,\qquad uvw+b=0,\qquad s^2+4t+4as+16c=0.
\]
Mit Hülfe der ersten dieser drei Gleichungen wird die dritte
\[
 t=a^2-4c,
\]
und, wenn man für $s,t$ die Werthe (2) zurücksetzt,
\begin{equation}
\begin{array}{rcl}
 u^2+v^2+w^2&=&-2a,\\
 v^2w^2+w^2u^2+u^2v^2&=&a^2-4c,\\
 uvw&=&-b.
\end{array}
\tag{3}
\end{equation}
Nach \S.\,7 sind diese Gleichungen dann und nur dann befriedigt, wenn $u^2,v^2,w^2$ die Wurzeln der cubischen Gleichung
\begin{equation}
 z^3+2az^2+(a^2-4c)z-b^2=0
\tag{4}
\end{equation}
sind, und wenn die Vorzeichen von $u,v,w$ so bestimmt werden, dass die letzte der Gleichungen (3) befriedigt ist. Diese letztere Bedingung lässt noch vier verschiedene Vorzeichenbestimmungen zu, so dass man die vier Wurzeln der biquadratischen Gleichung in folgender Weise erhält:
\begin{equation}
\begin{array}{rcl}
2\alpha&=&u+v+w,\\
2\beta&=&u-v-w,\\
2\gamma&=&-u+v-w,\\
2\delta&=&-u-v+w.
\end{array}
\tag{5}
\end{equation}
Die Lösung der biquadratischen Gleichung ist damit auf die der cubischen Gleichung (4) zurückgeführt.

Diese Gleichung heisst eine \gesperrt{cubische Resolvente} der biquadratischen Gleichung.
